\documentclass[12pt,a4paper]{article}

% === Packages ===
\usepackage[utf8]{inputenc}
\usepackage[T1]{fontenc}
\usepackage[brazil]{babel}
\usepackage{amsmath,amssymb,amsthm}
\usepackage{graphicx}
\usepackage{booktabs}
\usepackage{multirow}
\usepackage{natbib}
\usepackage{hyperref}
\usepackage[margin=2.5cm]{geometry}
\usepackage{setspace}
\onehalfspacing
\graphicspath{{figures/}}

% Refer\^encias cruzadas para o artigo principal
\usepackage{xr-hyper}
\externaldocument{main_pt}

% Renumera\c{c}\~ao com prefixo de letra
\renewcommand{\thesection}{\Alph{section}}
\renewcommand{\theequation}{A.\arabic{equation}}
\renewcommand{\thetable}{A.\arabic{table}}
\renewcommand{\thefigure}{A.\arabic{figure}}

% === Metadata ===
\title{Ap\^endice Online de\\
``Quanto de produtividade precisamos para reduzir a jornada de trabalho?''}
\author{Victor Rangel\thanks{Insper. E-mail: \texttt{victorr@insper.edu.br}.}}
\date{Mar\c{c}o de 2026}

\begin{document}

\maketitle

Este ap\^endice online re\'une fontes de dados, omiss\~oes do equil\'ibrio parcial, detalhes da calibra\c{c}\~ao, grades de sensibilidade, decomposi\c{c}\~ao setorial, vis\~ao por tipo do bem-estar, decomposi\c{c}\~ao do produto, tabela hist\'orica de PTF, deriva\c{c}\~oes e a robustez sob $e(h)$ assim\'etrico, todas referenciadas no texto principal. Numera\c{c}\~oes de se\c{c}\~oes, equa\c{c}\~oes, tabelas e figuras s\~ao prefixadas com ``A.''. Refer\^encias cruzadas ao artigo principal resolvem para os labels do documento principal.

\tableofcontents
\clearpage

\section{Fontes de dados}\label{app:data}

Todos os dados prov\^em de fontes p\'ublicas. PNAD Cont\'inua (IBGE/SIDRA), tabelas 4093, 4097, 6374 e 6413 (trimestrais, 2023--2025), fornece taxa de informalidade, posi\c{c}\~ao na ocupa\c{c}\~ao, horas m\'edias, distribui\c{c}\~ao por tamanho de firma e o peso CES formal $\omega=0{,}622$ ($1$ menos a informalidade restrita PNAD 2025). RAIS/CAGED (Minist\'erio do Trabalho) fornece registros administrativos anuais do emprego formal, usados para distribui\c{c}\~ao por tamanho de firma e parcelas formal/informal. A Penn World Table 10.01 fornece a parcela do trabalho para o Brasil ($\text{labsh}=0{,}578$), ajustada via \citet{Gollin2002} para $\alpha=0{,}35$. DIEESE (2024) fornece a distribui\c{c}\~ao de horas na for\c{c}a de trabalho formal, $\theta=(0{,}085,\,0{,}269,\,0{,}646)$ para os bins $\{36,\,40,\,44\}$. Um exerc\'icio de auditoria no pacote de replica\c{c}\~ao verifica que os momentos n\~ao targeteados no texto principal n\~ao s\~ao consequ\^encias alg\'ebricas da calibra\c{c}\~ao.

\clearpage
\section{O que o equil\'ibrio parcial omite}\label{app:pe_omits}

O arcabou\c{c}o est\'atico de equil\'ibrio parcial exclui sete canais de equil\'ibrio geral. A acumula\c{c}\~ao de capital responde \`a menor produtividade marginal do capital ap\'os a contra\c{c}\~ao do trabalho efetivo. Sal\'arios e pre\c{c}os ajustam-se em equil\'ibrio geral, restaurando parcialmente a demanda por trabalho formal. A participa\c{c}\~ao na for\c{c}a de trabalho pode subir com o lazer. Autoridades monet\'aria e fiscal podem acomodar. Trabalhadores podem realocar-se entre setores. A decis\~ao de formalidade pode mover-se em qualquer dire\c{c}\~ao. A partilha de trabalho \'e desligada pela hip\'otese de $N$ fixo. Todos os sete canais exceto o \'ultimo t\^em sinal esperado negativo sobre $A_{\text{req}}$.

Um bound textbook sobre o canal do capital d\'a a ordem de magnitude. Com elasticidade de custo de uso $\varepsilon_K\in[0{,}2,\,0{,}4]$ \citep{Caballero1999} e parcela de capital $\alpha=0{,}35$, em horizonte de $3$--$5$ anos o $A_{\text{req}}$ efetivo sob a especifica\c{c}\~ao conservadora cai de $8{,}18\%$ para $6{,}4\%$--$7{,}2\%$, e sob a preferida para $4{,}8\%$--$5{,}6\%$. Os seis canais restantes s\~ao pequenos individualmente e de sinal misto. Os n\'umeros centrais devem ser lidos como limites superiores de curto prazo sobre o ganho de PTF necess\'ario. O horizonte de $1$--$2$ anos do modelo n\~ao \'e o mesmo que os 30 anos da s\'erie PWT usada como \^ancora. A \^ancora \'e lida como o ganho cumulativo de PTF necess\'ario ao longo do per\'iodo de implementa\c{c}\~ao faseada de $3$--$10$ anos.

\clearpage
\section{Detalhes da calibra\c{c}\~ao}\label{app:calibration_details}

\subsection{Elasticidade horas--produto e forma da efici\^encia}

A evid\^encia de Pencavel documenta produtividade marginal decrescente de horas em jornadas longas para um setor industrial espec\'ifico. \citet{CollewetSauermann2017} encontram perdas an\'alogas em horas longas em um painel holand\^es de call-center. \citet{BickFuchsSchundelnLagakos2018} trazem evid\^encia transnacional consistente com produtividade marginal decrescente de horas, embora seu foco primordial seja a rela\c{c}\~ao entre horas e renda. Nenhum dos tr\^es trabalhos identifica uma localiza\c{c}\~ao precisa do pico nem uma resposta de efici\^encia abaixo do pico. A escolha de colocar o pico em $h^*=40$ e comparar uma especifica\c{c}\~ao preferida flat-below (sem penalidade para horas menores) com uma conservadora sim\'etrica (curvatura igual em ambos os lados) \'e uma escolha de modelagem informada por essa literatura, e n\~ao uma estimativa direta dela. A Se\c{c}\~ao~\ref{app:asymmetric} reporta quatro formas alternativas.

\subsection{Efici\^encia informal e custos de ajuste}

A efici\^encia informal $\eta_I=0{,}40$ \'e fixada no ponto m\'edio da raz\~ao salarial informal/formal LPS (faixa $1/3$ a $1/2$). Sob concorr\^encia perfeita e salarista, essa \'e a identifica\c{c}\~ao estrutural correta do par\^ametro de produtividade marginal relativa no CES, preferindo a raz\~ao salarial \`a raz\~ao VA/empregado, que \'e contaminada por gaps de capital por trabalhador. Os custos de ajuste $\gamma_F^S=0{,}12$ e $\gamma_F^L=0{,}03$ refletem o custo marginal efetivo de compliance CLT/FGTS para firmas pequenas e grandes. Sensibilidade em $\gamma_F^S\in[0{,}08,\,0{,}16]$ e $\gamma_F^L\in[0{,}01,\,0{,}05]$ desloca $A_{\text{req}}$ em menos de $\pm 0{,}2$ p.p. O custo convexo de informalidade $\pi_m$ \'e calibrado endogenamente para garantir cunhas n\~ao negativas no target de informalidade.

\subsection{Reconcilia\c{c}\~ao com elasticidades de substitui\c{c}\~ao maiores}\label{app:sigma_discussion}

Estudos de crise de refugiados reportam elasticidades de substitui\c{c}\~ao entre tipos de trabalhadores de $\sigma\approx 11$--$30$ \citep{OttavianoPeri2012, BaharDiTellaGulek2025}. Essas magnitudes n\~ao s\~ao diretamente compar\'aveis a $\sigma_{\text{sub}}$ como definido aqui por tr\^es raz\~oes. A margem difere. Estudos de refugiados medem substitui\c{c}\~ao entre trabalhadores de origens diferentes em equil\'ibrio, enquanto $\sigma_{\text{sub}}$ captura substitui\c{c}\~ao intrafirma entre regimes de emprego formal e informal. O horizonte difere. Estudos de refugiados identificam respostas de m\'edio prazo que incluem entrada de firmas e realoca\c{c}\~ao organizacional, enquanto este modelo mant\'em a distribui\c{c}\~ao de firmas fixa. A unidade difere. Estudos de refugiados incluem trabalhadores aut\^onomos, enquanto este agregador trata de emprego assalariado. Modelos DSGE de informalidade \citep{LeyvaUrrutia2020, McKiernan2021} tamb\'em usam par\^ametros $\sigma$ que governam mobilidade setorial em vez de substitui\c{c}\~ao intrafirma. Em \citet{Ulyssea2018}, trabalhadores formais e informais s\~ao substitutos perfeitos dentro da firma, com diferen\c{c}as de produtividade capturadas inteiramente pela distribui\c{c}\~ao de tipos de firmas. A calibra\c{c}\~ao $\sigma=1{,}326$ pode ser lida como uma decomposi\c{c}\~ao param\'etrica que atribui parte do gap de produtividade formal--informal a $\eta_I=0{,}40$ e parte a $\sigma>1$ (substitutos imperfeitos mas n\~ao complementos), ambos disciplinados pelo pr\^emio salarial observado.

\subsection{Disciplinando $\sigma_{\text{sub}}$: resultados completos}

Para cada $\sigma$ em grade $[0{,}40,\,3{,}00]$, recalibra-se o modelo completo e calcula-se o pr\^emio salarial formal--informal implicado pela CPO do CES,
\[
R(\sigma)=\frac{\omega}{1-\omega}\cdot\left(\frac{L_F}{L_I}\right)^{\rho-1}\cdot z,
\]
onde $z=\text{eff}_{hF}/(\eta_I h_I e(h_I))$. O mapeamento \'e mon\'otono decrescente em $\sigma$. Fixando $R\in[1{,}15,\,1{,}55]$ obt\'em-se $\sigma\in[1{,}116,\,1{,}469]$ e o $A_{\text{req}}\in[7{,}60\%,\,8{,}70\%]$ conservador. Corbi-Ferraz-Narita (2024) e \citet{DerenoncourtGerardLagosMontialoux2025} encontram elasticidades de realoca\c{c}\~ao pequenas consistentes com essa faixa.

\clearpage
\section{Sensibilidade}\label{app:sensitivity}

A Tabela~\ref{tab:sensitivity_univariate} reporta $A_{\text{req}}$ sob varia\c{c}\~ao univariada de cada par\^ametro dentro de faixas economicamente disciplinadas. Tr\^es par\^ametros carregam o peso, $\sigma$, $\omega$ e $\eta_I$. Os demais primitivos deslocam $A_{\text{req}}$ em no m\'aximo $\pm 0{,}3$ p.p. A Tabela~\ref{tab:heatmap} reporta a grade completa $(\sigma,\omega)$. O resultado qualitativo de que $A_{\text{req}}$ excede $4\%$ em todas as combina\c{c}\~oes da grade \'e plenamente robusto.

\begin{table}[htbp]
\centering
\caption{Sensibilidade univariada dentro de faixas economicamente disciplinadas (conservadora).}\label{tab:sensitivity_univariate}
\small
\begin{tabular}{llrl}
\hline\hline
Par\^ametro & Faixa disciplinada & Faixa de $A_{\text{req}}$ & Hierarquia \\
\hline
$\sigma_{\text{sub}}$ & $[1{,}12,\,1{,}47]$ & $[7{,}60\%,\,8{,}70\%]$ & Alta (load-bearing) \\
$\omega$ & $[0{,}58,\,0{,}66]$ & $[7{,}45\%,\,8{,}72\%]$ & Alta \\
$\eta_I$ & $[0{,}33,\,0{,}50]$ & $[7{,}55\%,\,8{,}55\%]$ & Moderada \\
$\alpha$ & $[0{,}30,\,0{,}40]$ & $[7{,}55\%,\,8{,}85\%]$ & Moderada \\
$e_q$ & $[0{,}40,\,0{,}80]$ & $[7{,}95\%,\,8{,}41\%]$ & Baixa \\
$h_I$ & $[38,\,46]$ & $[8{,}00\%,\,8{,}24\%]$ & Baixa \\
$\nu$ & $[1{,}0,\,3{,}0]$ & $\Delta$CV: $[-3{,}80\%,\,-3{,}55\%]$ & Baixa (apenas bem-estar) \\
\hline\hline
\end{tabular}
\\[2pt]
\begin{minipage}{0.95\textwidth}
{\footnotesize \textit{Notas:} Cada linha varia um par\^ametro dentro de sua faixa consistente com os dados e mant\'em os demais no baseline. ``Faixa disciplinada'' \'e o intervalo derivado de evid\^encia externa (Se\c{c}\~ao~\ref{app:calibration_details}). ``Faixa de $A_{\text{req}}$'' \'e o intervalo implicado para o ganho requerido de PTF sob a especifica\c{c}\~ao conservadora de $e(h)$ no teto $44\to 36$h. ``Hierarquia'' classifica cada par\^ametro pela amplitude de seu intervalo de $A_{\text{req}}$.}
\end{minipage}
\end{table}

\begin{table}[htbp]
\centering
\caption{$A_{\text{req}}$ (\%) por $\sigma$ e $\omega$ em $\eta_I=0{,}40$ (conservadora).}\label{tab:heatmap}
\small
\begin{tabular}{r|rrr}
\hline\hline
$\sigma$ \textbackslash{} $\omega$ & 0{,}58 & 0{,}622 & 0{,}66 \\
\hline
1{,}116 & 7{,}00 & 7{,}50 & 7{,}96 \\
\textbf{1{,}326} & \textbf{7{,}65} & \textbf{8{,}18} & \textbf{8{,}65} \\
1{,}469 & 8{,}03 & 8{,}58 & 9{,}06 \\
\hline\hline
\end{tabular}
\\[2pt]
\begin{minipage}{0.80\textwidth}
{\footnotesize \textit{Notas:} Cada c\'elula \'e o ganho requerido de PTF (\%) no teto $44\to 36$h sob a especifica\c{c}\~ao conservadora sim\'etrica de $e(h)$, para o par $(\sigma,\omega)$ indicado e com os demais par\^ametros no baseline. A linha em negrito $\sigma=1{,}326$ \'e o valor central calibrado. O intervalo consistente com MNR \'e $[1{,}116,\,1{,}469]$.}
\end{minipage}
\end{table}

\clearpage
\section{Decomposi\c{c}\~ao setorial}\label{app:sectoral}

Esta se\c{c}\~ao aplica o arcabou\c{c}o nacional separadamente a tr\^es setores amplos (agricultura, ind\'ustria, servi\c{c}os) com par\^ametros setoriais obtidos dos microdados da PNAD Cont\'inua Q4 2024 (209{,}311 pessoas ocupadas, ponderadas pela amostra). Par\^ametros estruturais comuns ($\alpha$, $\sigma_{\text{sub}}$, $\omega$, $\kappa$, $h^*$) s\~ao preservados da calibra\c{c}\~ao nacional. Cada setor \'e tratado como inst\^ancia independente do modelo nacional. O exerc\'icio n\~ao modela mobilidade intersetorial ou rela\c{c}\~oes insumo-produto.

A Tabela~\ref{tab:sectoral_params} reporta par\^ametros setoriais e as cunhas calibradas. A informalidade da agricultura \'e de $71{,}8\%$, quase o dobro da m\'edia nacional de $44{,}2\%$ sob a defini\c{c}\~ao ampla do IBGE. Ind\'ustria ($44{,}4\%$) e servi\c{c}os ($41{,}2\%$) s\~ao mais pr\'oximos da m\'edia. As cunhas calibradas variam de $\tau_s=1{,}48$ em servi\c{c}os a $\tau_s=5{,}63$ na agricultura. A cunha agr\'icola \'e comparada com as obriga\c{c}\~oes patronais FUNRURAL-SENAR-INSS de $30$--$35\%$ dos sal\'arios mais custos de enforcement ligados \`a sazonalidade \citep{LemosGalvao2020}. A magnitude \'e grande e n\~ao implaus\'ivel dada a fric\c{c}\~ao institucional alta sobre uma base formal pequena.

\begin{table}[htbp]
\centering
\caption{Par\^ametros setoriais e cunhas calibradas.}\label{tab:sectoral_params}
\small
\begin{tabular}{lcccccccr}
\hline\hline
Setor & $\lambda_s$ & VAB & Inf.\ & $\theta_{36}$ & $\theta_{40}$ & $\theta_{44}$ & Hrs (F) & $\tau_s$ \\
\hline
Agricultura & 0{,}076 & 0{,}077 & 71{,}8\% & 0{,}101 & 0{,}269 & 0{,}630 & 44{,}1 & 5{,}63 \\
Ind\'ustria & 0{,}205 & 0{,}259 & 44{,}4\% & 0{,}048 & 0{,}417 & 0{,}535 & 42{,}5 &  2{,}10 \\
Servi\c{c}os & 0{,}720 & 0{,}665 & 41{,}2\% & 0{,}171 & 0{,}410 & 0{,}419 & 40{,}8 &  1{,}48 \\
\hline
Nacional    & 1{,}000 & 1{,}000 & 44{,}2\% & 0{,}143 & 0{,}406 & 0{,}451 & 41{,}2 & --- \\
\hline\hline
\end{tabular}
\\[2pt]
{\footnotesize $\lambda_s$: parcela de emprego, PNAD Q4 2024. VAB: parcela do valor adicionado, IBGE Contas Regionais 2021. $\theta_b$: parcela de trabalhadores formais no bin $b$. Horas: V4039 (habitual). Informalidade: defini\c{c}\~ao ampla (VD4009 + V4017 CNPJ).}
\end{table}

A Tabela~\ref{tab:sectoral_results} reporta os principais resultados setoriais. O $A_{\text{req}}$ agregado imp\~oe um multiplicador de PTF uniforme em todos os tr\^es setores e resolve por bisse\c{c}\~ao. Sob a reforma completa, a agricultura requer $A_{\text{req}}=7{,}60\%$, a ind\'ustria $8{,}05\%$ e servi\c{c}os $6{,}91\%$, com agregado de $7{,}21\%$. O $A_{\text{req}}$ da ind\'ustria supera ligeiramente o da agricultura. Com $\sigma>1$ e $\eta_I=0{,}40$, a fric\c{c}\~ao de realoca\c{c}\~ao escala com o tamanho da base formal, que \'e maior na ind\'ustria do que na agricultura. Todos os tr\^es setores v\^eem a informalidade subir em resposta \`a reforma (agricultura $+2{,}88$ p.p., ind\'ustria $+2{,}73$ p.p., servi\c{c}os $+2{,}10$ p.p.), consistente com trabalho formal--informal moderadamente substitu\'ivel.

\begin{table}[htbp]
\centering
\caption{Resultados setoriais: ganho requerido de PTF e exposi\c{c}\~ao dos trabalhadores.}\label{tab:sectoral_results}
\small
\begin{tabular}{llrrrrc}
\hline\hline
Reforma & Setor & $A_{\text{req}}$ & $\Delta Y$ & $\Delta$Inf & Afetados & Raz\~ao \\
 & & (\%) & (\%) & (p.p.) & (\% do setor) & ($A_{\text{req}}$/Af.) \\
\hline
\multirow{4}{*}{$44\to 40$h}
& Agricultura & 1{,}8 & $-$2{,}4 & $+$0{,}6 & 17{,}8 & 10{,}1 \\
& Ind\'ustria & 1{,}6 & $-$1{,}6 & $+$0{,}4 & 29{,}7 & 5{,}4 \\
& Servi\c{c}os & 1{,}3 & $-$1{,}3 & $+$0{,}4 & 24{,}7 & 5{,}2 \\
& \textit{Agregado} & \textit{1{,}4} & \textit{$-$1{,}4} & \textit{$+$0{,}4} & \textit{25{,}2} & \\
\hline
\multirow{4}{*}{$44\to 36$h}
& Agricultura & 7{,}6 & $-$9{,}4 & $+$2{,}9 & 25{,}3 & 30{,}0 \\
& Ind\'ustria & 8{,}0 & $-$7{,}9 & $+$2{,}7 & 52{,}9 & 15{,}2 \\
& Servi\c{c}os & 6{,}9 & $-$6{,}7 & $+$2{,}1 & 48{,}7 & 14{,}2 \\
& \textit{Agregado} & \textit{7{,}2} & \textit{$-$7{,}2} & \textit{$+$2{,}3} & \textit{47{,}8} & \\
\hline\hline
\end{tabular}
\end{table}

\begin{figure}[htbp]
\centering
\includegraphics[width=0.92\textwidth]{fig_sectoral_areq.pdf}
\caption{Ganho requerido de PTF $A_{\text{req}}$ por setor sob as reformas moderada $44\to 40$h e completa $44\to 36$h, calibra\c{c}\~ao conservadora.}\label{fig:sectoral_areq}
\vspace{2pt}
\begin{minipage}{0.92\textwidth}
{\footnotesize \textit{Notas:} Valores calculados separadamente para cada setor sob a especifica\c{c}\~ao conservadora sim\'etrica de $e(h)$. Par\^ametros espec\'ificos por setor (informalidade, distribui\c{c}\~ao de horas, cunha $\tau_s$) v\^em da PNAD Cont\'inua Q4 2024. Os par\^ametros estruturais comuns ($\alpha$, $\sigma_{\text{sub}}$, $\omega$, $\kappa$, $h^*$) s\~ao mantidos nos valores nacionais.}
\end{minipage}
\end{figure}

Os servi\c{c}os respondem por $71{,}0\%$ do produto de base implicado pelo modelo e contribuem com $4{,}8$ dos $7{,}2$ p.p.\ de perda agregada de produto. Uma pol\'itica setorial que isentasse a agricultura do teto reduziria o $A_{\text{req}}$ agregado em menos de $0{,}2$ p.p. A decomposi\c{c}\~ao trata cada setor como inst\^ancia independente do modelo nacional e abstrai da realoca\c{c}\~ao intersetorial de trabalho e capital. Tamb\'em colapsa a heterogeneidade por tamanho de firma em uma \'unica firma representativa por setor. Um arcabou\c{c}o multissetorial mais rico com mobilidade de fatores \'e uma extens\~ao natural.

\clearpage
\section{Bem-estar: vis\~ao por tipo}\label{app:welfare_additional}

A Tabela~\ref{tab:welfare_conditional} reporta uma decomposi\c{c}\~ao por tipo do efeito de bem-estar na reforma $44\to 36$h sob a especifica\c{c}\~ao conservadora, usando parcelas da folha de pagamento CES. Trabalhadores que permanecem formais trabalham aproximadamente 6 horas semanais a menos, com perda de consumo em torno de $6\%$ e ganho l\'iquido de bem-estar sob prefer\^encias GHH com $\nu=2$. O bem-estar dos trabalhadores informais \'e pouco afetado porque $h_I=44$ n\~ao est\'a sujeito ao teto. A massa de trabalhadores desformalizados na margem enfrenta perdas acentuadas. Sob o moderado $\eta_I=0{,}40$, essa massa \'e material e menor do que sob especifica\c{c}\~oes de menor efici\^encia. Os n\'umeros por tipo s\~ao uma aproxima\c{c}\~ao de primeira ordem. O CV exato do agente representativo na Equa\c{c}\~ao~\eqref{eq:cv} do artigo principal calcula o composto em consumo e horas agregados. O gap entre a soma ponderada por tipo e o agregado \'e a corre\c{c}\~ao tipo-Jensen implicada pela concavidade do composto GHH.

\begin{table}[htbp]
\centering
\caption{Bem-estar por tipo de trabalhador ($44\to 36$h, conservadora).}\label{tab:welfare_conditional}
\small
\begin{tabular}{lcc}
\hline\hline
Tipo de trabalhador & $\Delta$CV$_k$ (primeira ordem) & Parcela $s_k$ \\
\hline
Formais que permanecem & $+2{,}0\%$ & $60{,}5\%$ \\
Informais (inalterados) & $\approx 0$ & $37{,}5\%$ \\
Desformalizados & $-30{,}0\%$ & $2{,}0\%$ \\
\hline
Agregado, Equa\c{c}\~ao~\eqref{eq:cv} do principal & $-3{,}63\%$ & $100\%$ \\
\hline\hline
\end{tabular}
\\[2pt]
\begin{minipage}{0.90\textwidth}
{\footnotesize \textit{Notas:} $\Delta$CV$_k$ \'e a varia\c{c}\~ao compensat\'oria de primeira ordem aplicada ao tipo $k$ (GHH, $\nu=2$), avaliada no consumo e nas horas pr\'oprios de cada tipo. $s_k$ \'e a parcela da folha de pagamento CES usada como peso de agrega\c{c}\~ao. O valor do agente representativo pela Equa\c{c}\~ao~\eqref{eq:cv} do artigo principal \'e $-3{,}63\%$. O gap entre a soma ponderada por tipo e o valor agregado \'e a corre\c{c}\~ao tipo-Jensen implicada pela concavidade do composto GHH.}
\end{minipage}
\end{table}

A Figura~\ref{fig:welfare_threshold} mostra a resposta de bem-estar como fun\c{c}\~ao de um ganho ex\'ogeno de PTF no teto de 36h sob as duas especifica\c{c}\~oes. Os limiares de neutralidade de bem-estar s\~ao gerados pelo pipeline de replica\c{c}\~ao e marcados na figura.

\begin{figure}[htbp]
\centering
\includegraphics[width=0.85\textwidth]{slide10c_welfare_threshold.png}
\caption{Bem-estar $\Delta$CV como fun\c{c}\~ao do ganho de PTF no teto de 36h.}\label{fig:welfare_threshold}
\vspace{2pt}
\begin{minipage}{0.85\textwidth}
{\footnotesize \textit{Notas:} Eixo horizontal: ganho ex\'ogeno \'unico de PTF aplicado sobre a reforma $44\to 36$h. Eixo vertical: bem-estar agregado $\Delta$CV da Equa\c{c}\~ao~\eqref{eq:cv} do artigo principal. As curvas e os limiares s\~ao gerados por \texttt{src/tables\_figures/plot\_welfare\_schedule.py}.}
\end{minipage}
\end{figure}

\begin{table}[htbp]
\centering
\caption{Schedule de bem-estar, especifica\c{c}\~ao conservadora.}\label{tab:welfare_schedule}
\small
\begin{tabular}{rrrrl}
\hline\hline
Teto de horas & $A_{\text{req}}$ & $\Delta Y$ & $\Delta$CV & $\Delta$Inf \\
\hline
44 & 0{,}00\% & 0{,}00\% & 0{,}00\% & 0{,}00 p.p. \\
43 & 0{,}26\% & $-$0{,}27\% & $+$0{,}42\% & $+$0{,}07 p.p. \\
42 & 0{,}66\% & $-$0{,}70\% & $+$0{,}63\% & $+$0{,}16 p.p. \\
41 & 1{,}22\% & $-$1{,}27\% & $+$0{,}63\% & $+$0{,}29 p.p. \\
40 & 1{,}92\% & $-$2{,}00\% & $+$0{,}41\% & $+$0{,}46 p.p. \\
39 & 3{,}14\% & $-$3{,}23\% & $-$0{,}18\% & $+$0{,}75 p.p. \\
38 & 4{,}58\% & $-$4{,}65\% & $-$1{,}06\% & $+$1{,}09 p.p. \\
37 & 6{,}26\% & $-$6{,}25\% & $-$2{,}21\% & $+$1{,}48 p.p. \\
\textbf{36} & \textbf{8{,}18\%} & $-$\textbf{8{,}02\%} & $-$\textbf{3{,}63\%} & $+$\textbf{1{,}92 p.p.} \\
35 & 10{,}57\% & $-$10{,}14\% & $-$5{,}45\% & $+$2{,}46 p.p. \\
\hline\hline
\end{tabular}
\\[2pt]
\begin{minipage}{0.85\textwidth}
{\footnotesize \textit{Notas:} Todas as linhas s\~ao resultados do modelo sob a especifica\c{c}\~ao conservadora sim\'etrica de $e(h)$, medidos em rela\c{c}\~ao ao baseline de 44h. $A_{\text{req}}$ \'e o ganho requerido de PTF para o teto indicado. $\Delta Y$ \'e a varia\c{c}\~ao agregada do produto. $\Delta$CV \'e o bem-estar agregado pela Equa\c{c}\~ao~\eqref{eq:cv} do artigo principal. $\Delta$Inf \'e a varia\c{c}\~ao da informalidade agregada. A linha em negrito \'e o teto proposto de 36h.}
\end{minipage}
\end{table}

\clearpage
\section{Decomposi\c{c}\~ao do produto}\label{app:decomposition}

\begin{figure}[htbp]
\centering
\includegraphics[width=0.92\textwidth]{fig_decomposition.pdf}
\caption{Decomposi\c{c}\~ao de $\Delta Y$ no teto $44\to 36$h sob cada calibra\c{c}\~ao.}\label{fig:decomposition}
\vspace{2pt}
\begin{minipage}{0.92\textwidth}
{\footnotesize \textit{Notas:} As barras s\~ao geradas por \texttt{src/tables\_figures/plot\_main\_figures.py}. O canal mec\^anico de horas \'e o efeito direto do teto menor de horas formais, mantendo fixa a efici\^encia por hora. O canal de efici\^encia \'e a mudan\c{c}a na efici\^encia hor\'aria formal $e(h)$. O res\'iduo inclui a resposta de realoca\c{c}\~ao formal--informal e o termo cruzado CES. Os valores exatos s\~ao gravados em \texttt{output/validation/decomposition\_results.json}.}
\end{minipage}
\end{figure}

\clearpage
\section{Tabela hist\'orica de PTF}\label{app:tfp_history_table}

\begin{table}[htbp]
\centering
\caption{Crescimento hist\'orico da PTF brasileira e anos necess\'arios para acumular $A_{\text{req}}$.}\label{tab:tfp_history}
\small
\begin{tabular}{lrrrr}
\hline\hline
Horizonte & Dura\c{c}\~ao (anos) & Crescim.\ PTF (\%/ano) & Anos at\'e $A_{\text{req}}$ & Viabilidade \\
\hline
1954--2019 & 65 & $+0{,}29$ & 22{,}4 & vi\'avel \\
1954--1980 & 26 & $+2{,}29$ & 2{,}8 & vi\'avel \\
1980--2000 & 20 & $-1{,}31$ & $\infty$ & \textit{nunca} \\
1990--2019 & 29 & $-0{,}60$ & $\infty$ & \textit{nunca} \\
1995--2019 & 24 & $-0{,}83$ & $\infty$ & \textit{nunca} \\
2000--2019 & 19 & $-0{,}72$ & $\infty$ & \textit{nunca} \\
2003--2013 & 10 & $+0{,}26$ & 24{,}7 & vi\'avel \\
2010--2019 & 9 & $-1{,}88$ & $\infty$ & \textit{nunca} \\
\hline\hline
\end{tabular}
\\[2pt]
\begin{minipage}{0.95\textwidth}
{\footnotesize \textit{Notas:} Crescimento anualizado \'e $(\text{PTF}_1/\text{PTF}_0)^{1/T}-1$. A coluna ``Anos at\'e $A_{\text{req}}$'' usa o baseline preferido $A_{\text{req}}=6{,}63\%$ e \'e calculada como $\ln(1+A_{\text{req}})/\ln(1+g)$. Horizontes com crescimento n\~ao positivo s\~ao reportados como ``nunca''. \textit{Fonte:} Penn World Table 10.01, vari\'avel \texttt{rtfpna} (PTF real a pre\c{c}os nacionais constantes).}
\end{minipage}
\end{table}

\clearpage
\section{Deriva\c{c}\~oes}\label{app:derivations}

\subsection{$\kappa$ a partir de $e_q$}\label{app:kappa_deriv}

Com $e(h)=\exp\{-\kappa(h-h^*)^2\}$, a elasticidade horas--produto em $h_{\text{ref}}$ \'e
\[
e_q = \frac{d\ln(h\cdot e(h))}{d\ln h}\bigg|_{h=h_{\text{ref}}} = 1 - 2\kappa\,h_{\text{ref}}(h_{\text{ref}}-h^*),
\]
de modo que
\begin{equation}\label{eq:kappa_derivation}
\kappa=\frac{1-e_q}{2 h_{\text{ref}}(h_{\text{ref}}-h^*)}.
\end{equation}
Com $e_q=0{,}60$, $h_{\text{ref}}=42{,}244$ e $h^*=40$, $\kappa=2{,}110\times 10^{-3}$, ent\~ao $e(36)=e(44)=0{,}9668$ e $e(40)=1$.

\subsection{Condi\c{c}\~ao de primeira ordem da firma}

O problema da firma tem
\begin{equation}\label{eq:foc_NF}
\frac{\partial Y_g}{\partial L_{F,g}}\cdot h_{F,\text{avg}}\cdot e_F\cdot\omega\cdot\left(\frac{L_{F,g}}{L_g}\right)^{\rho-1}=\tau_g+\gamma_{F,g}(N_{F,g}-\bar{N}_{F,g})+\pi_m N_{I,g},
\end{equation}
onde $\partial Y_g/\partial L_{F,g}=(1-\alpha)Y_g/L_g\cdot(\partial L_g/\partial L_{F,g})$ e $\partial L_g/\partial L_{F,g}=\omega(L_{F,g}/L_g)^{\rho-1}$.

\subsection{Monotonicidade da bisse\c{c}\~ao de $A_{\text{req}}$}

$A_{\text{req}}$ resolve a Equa\c{c}\~ao~\eqref{eq:areq} do artigo principal. A estrutura Cobb-Douglas $Y_g=A K_g^\alpha L_g^{1-\alpha}$ \'e linear em $A$ e o $L_g(A)$ \'otimo \'e fracamente crescente em $A$, ent\~ao $\partial\sum_g Y_g/\partial A>0$ e a bisse\c{c}\~ao converge para um \'unico $A_{\text{req}}$.

\clearpage
\section{Robustez do $e(h)$ assim\'etrico}\label{app:asymmetric}

A Tabela~\ref{tab:asymmetric_eh} reporta $A_{\text{req}}$ sob quatro especifica\c{c}\~oes para $e(h)$. O $e(h)=\exp\{-\kappa(h-40)^2\}$ sim\'etrico de baseline com $\kappa=2{,}11\times 10^{-3}$ \'e comparado com tr\^es alternativas assim\'etricas que variam a curvatura abaixo de $h^*=40$ enquanto mant\^em a curvatura acima de $h^*$ em $\kappa$.

\begin{table}[htbp]
\centering
\caption{$A_{\text{req}}$ sob $e(h)$ assim\'etrico.}\label{tab:asymmetric_eh}
\small
\begin{tabular}{llrr}
\hline\hline
Cen\'ario & Hip\'otese abaixo de $h^*$ & $e(36)$ & $A_{\text{req}}$ \\
\hline
Sim\'etrico (conservador) & $\kappa_{\text{abaixo}}=\kappa$ & 0{,}967 & 8{,}18\% \\
Meia-curvatura abaixo & $\kappa_{\text{abaixo}}=0{,}5\kappa$ & 0{,}983 & 7{,}40\% \\
Curvatura dupla abaixo & $\kappa_{\text{abaixo}}=2\kappa$ & 0{,}935 & 9{,}73\% \\
Flat abaixo (preferida) & $\kappa_{\text{abaixo}}=0$ & 1{,}000 & 6{,}63\% \\
\hline\hline
\end{tabular}
\\[2pt]
\begin{minipage}{0.90\textwidth}
{\footnotesize \textit{Notas:} A curvatura acima de $h^*=40$ \'e fixada em $\kappa=2{,}11\times 10^{-3}$ nos quatro cen\'arios. Apenas a curvatura abaixo do pico varia. $e(36)$ \'e a efici\^encia por hora no teto proposto de 36h (quanto mais distante de 1, maior a penalidade de efici\^encia). $A_{\text{req}}$ \'e o ganho requerido de PTF no teto $44\to 36$h sob cada cen\'ario.}
\end{minipage}
\end{table}

% === Bibliografia ===
\clearpage
\bibliographystyle{aer}
\bibliography{bibliography_clean}

\end{document}
